\section{The gap and the model}
\label{sec:intro}
Laughlin's wave function gave an explicit description of the fractional
quantum Hall state at filling $1/3$~\cite{Laughlin}. Haldane's spherical
geometry and pseudopotentials identify a particularly simple parent
Hamiltonian: every pair pays energy one in relative angular momentum one
and zero in the other pair sectors~\cite{Haldane}. Its exact Laughlin
zero mode is known, but a uniform lower bound on the positive spectrum
requires control of all competing many-particle states. We prove this
bound for the full interaction on the round sphere, resolving positively
the spherical fermionic Laughlin spectral-gap conjecture.

The gap question has accompanied this model since its early development.
Haldane and Rezayi's finite-size study supplied numerical evidence and
explicitly distinguished the expected hard-core gap from a formal
proof~\cite{HaldaneRezayi1985FiniteSize}. The density-wave variational
picture of Girvin, MacDonald, and Platzman explains important features of
the excitation spectrum~\cite{GMP}. A variational excitation energy is
an upper bound; the lower bound sought here must hold for every state
orthogonal to the zero modes.

\subsection{Normalization and main results}
For an integer \(Q\ge1\), let \(U_Q\) be the spin-\(Q/2\) representation of
\(SU(2)\), equipped with its invariant inner product. It is the one-particle
lowest Landau level on the sphere with \(Q\) flux quanta. Equivalently, it is
the space of homogeneous polynomials of degree \(Q\) in spinor coordinates
\((u,v)\). For \(2\le N\le Q+1\), define
\begin{equation}\label{eq:physicalH}
 H_{N,Q}=\left.\sum_{1\le i<j\le N}P_{ij}^{(1)}
          \right|_{\bigwedge^N U_Q},
\end{equation}
where \(P_{ij}^{(1)}\) is the orthogonal projector onto pair spin \(Q-1\)
in tensor positions \(i,j\). This is relative angular momentum one. The sum
preserves the antisymmetric subspace. Each unordered pair has coefficient
one, with no rescaling depending on \(N\) or \(Q\).

Set \(H_{0,Q}=H_{1,Q}=0\), and write
\[
 \FF_Q=\bigoplus_{N=0}^{Q+1}\bigwedge^N U_Q,\qquad
 H_Q=\bigoplus_{N=0}^{Q+1}H_{N,Q},\qquad
 \gamma_*=\frac{4616733319001}{10^{14}}> \frac1{25}.
\]
Our principal estimate concerns this entire Fock space.

\begin{theorem}\label{thm:fock}
For every \(0<\gamma<\gamma_*\), there is an integer \(Q_\gamma\) such that
for every integer \(Q\ge Q_\gamma\),
\begin{equation}\label{eq:architecture}
 H_Q^2\ge\gamma H_Q\qquad\text{on }\FF_Q.
\end{equation}
The threshold \(Q_\gamma\) is independent of particle number.
\end{theorem}

Thus the spectrum is contained in \(\{0\}\cup[\gamma,\infty)\).
In a particle sector without a zero mode, this statement bounds the energy
away from zero; it does not assert a gap between that sector's two lowest
eigenvalues. At Laughlin filling the zero mode is unique, and we obtain the
usual ground-state spectral gap.

\begin{corollary}\label{thm:main}
There is an integer \(N_0\ge2\) such that, for every \(N\ge N_0\) and
\(Q=3(N-1)\),
\begin{equation}\label{eq:main}
 H_{N,Q}\ge\frac1{25}\bigl(I-P_{\mathrm L,N}\bigr)
 \qquad\text{on }\bigwedge^N U_Q,
\end{equation}
where \(P_{\mathrm L,N}\) projects onto the line spanned by
\[
 \Psi_{\mathrm L,N}=\prod_{i<j}(u_iv_j-u_jv_i)^3.
\]
\end{corollary}

At fixed magnetic length, the sphere radius is proportional to \(\sqrt Q\),
so this is the two-dimensional thermodynamic limit. The thresholds in
Theorem~\ref{thm:fock} and Corollary~\ref{thm:main} are existential.
The proof does not give an explicit flux threshold or the endpoint
coefficient \(\gamma_*\) on finite spheres.

\subsection{Prior work}
The exact ground-state structure is an established part of the theory.
Besides Haldane's spherical construction, Trugman and Kivelson obtained
exact Laughlin ground states for repulsive interactions in a vanishing-range
limit~\cite{TrugmanKivelson1985}. The second-quantized algebra of
pseudopotentials was developed systematically by Ortiz, Nussinov, Dukelsky,
and Seidel~\cite{OrtizNussinovDukelskySeidel}. Chen and Seidel, and
Mazaheri, Ortiz, Nussinov, and Seidel, gave algebraic constructions and
characterizations of their zero modes~\cite{ChenSeidel2015,MazaheriOrtizNussinovSeidel2015}.
These results identify the zero modes that a spectral estimate must preserve. We include a short self-contained proof of the known spherical
zero-mode uniqueness at the Laughlin flux in Lemma~\ref{lem:kernel}.

For Laughlin states on a cylinder of any fixed radius, Jansen established
thermodynamic correlation limits and axial periodicity~\cite{Jansen2012}.
Exact ground states of related finite-range orbital chains, including
squeezed and matrix-product states, were constructed by Nakamura, Wang, and
Bergholtz~\cite{NakamuraWangBergholtz2012} and by Wang and
Nakamura~\cite{WangNakamura2013}. Nachtergaele, Warzel, and Young proved
uniform gaps for a truncated fermionic pseudopotential~\cite{NWY};
Warzel and Young obtained thin-torus bounds that separate bulk behavior
from open-chain edge modes~\cite{WY}. Their invariant-subspace method
also treats a truncated bosonic pseudopotential~\cite{WarzelYoung2023BulkEdge}.
The fermionic estimates already hold uniformly across particle sectors.
The distinction in Theorem~\ref{thm:fock} is the full spherical pair
interaction, with its fixed projector normalization and sphere radius
growing in the thermodynamic limit. The finite support used in our
auxiliary comparison operators does not truncate that interaction.

Recent work gives further structural and spectral information.
Lemm, Nachtergaele, Warzel, and Young derive recursions between particle
sectors and compare charge and neutral gaps~\cite{LemmNachtergaeleWarzelYoung2026Recursive};
a uniform bound for the full pseudopotential still requires control of
the model-specific ground-space overlap matrices in those criteria.
Schraven and Warzel prove entanglement-gap and clustering results for
Laughlin states on sufficiently thin cylinders~\cite{SchravenWarzel2026IMPS}.
The entanglement spectrum is distinct from the Hamiltonian spectrum
considered here. Density rigidity within the family of fully correlated
Laughlin states is another complementary form of incompressibility,
established by Lieb, Rougerie, and Yngvason~\cite{LRYRigidity}.
No density-screening assumption enters our argument.

Rougerie's survey records the spectral-gap conjecture and its geometric
formulations~\cite[Appendix~A]{Rougerie}. For $N$ fermions in the planar
lowest Landau level, consider the sum over unordered pairs of the orthogonal
projectors onto relative polynomial degree one, each with coefficient one
and with center-of-mass degree unrestricted. Section~\ref{sec:plane} proves
that its positive spectrum is bounded below by the endpoint $\gamma_*$,
uniformly in particle number $N\ge2$ and total polynomial degree $L$.
In particular, this proves the fermionic cubic case of the folklore
spectral-gap conjecture in \cite[Conjecture~A.1]{Rougerie}, whose degree range
is $L\le3N(N-1)/2$.

\subsection{Proof strategy}
Proving positivity of $H^2-\gamma H$ is a standard route to a gap,
underlying finite-size criteria~\cite{Knabe} and methods based on positive
local decompositions~\cite{Cruz,Rai}. Our argument uses rotations and
angular-momentum decomposition to construct such a comparison.
It starts from the normal-ordered expansion of \(H_Q^2\).
It has two-, three-, and four-body terms. We compare them with the
corresponding terms of a positive operator
\[
 \mathcal K_Q=(2Q-1)\int_{SU(2)}
                  \sum_{\rho=1}^{7}F_\rho(g)^\dagger F_\rho(g)\,dg,
\]
where each \(F_\rho\) is a specified linear combination of a pair
annihilator and operators that annihilate three particles and create one.
Here \(F_\rho(g)\) denotes rotation by \(g\), and Haar measure has mass
one. These auxiliary operators involve finitely many orbital modes;
the Hamiltonian itself is never truncated.

Rotational symmetry reduces the comparison to fixed spin blocks.
The three-body blocks have an explicit spectrum and a summable tail.
For four particles, coupling the interacting pair to the different spin
sectors of a second pair can produce several copies of the same total spin.
Their four-particle images under exterior multiplication can be linearly
dependent.
For the retained blocks, seven rows of rational data establish finite
inequalities in their large-flux limits. The remaining issue is to transfer
those inequalities to finite spheres without an error that grows with the
number of particles.

The key device is an exact diagonal change of variables for the pair
coefficients. On a fixed set of orbital modes, it converts finite-flux
annihilators into their limiting forms. Every perturbation is then bounded
by the local pair energy. Rotation averaging turns this into
\(o(1)H_Q\) on the full Fock space. This relative estimate remains valid
when the Gram matrices of these limiting four-particle images have kernels.
A second useful observation
is to apply Bessel's inequality to all equivalent spin copies at once;
this substantially reduces the four-body comparison error.

The comparisons leave a positive margin:
\begin{equation}\label{eq:budget}
 H_Q^2\ge\mathcal K_Q+
       \bigl(\gamma_*-o(1)\bigr)H_Q.
\end{equation}
All errors and convergence
thresholds are independent of particle number.
Figure~\ref{fig:transfer} isolates the passage responsible for this
uniformity. A small norm error on a growing many-particle space would not
suffice: it must vanish on the zero modes and be controlled by the pair
energy. The fixed-mode comparison gives precisely that control before
rotations distribute it over the sphere.

\begin{figure}[t]
\centering
\begin{tikzpicture}[
  box/.style={draw=blue!45!black,rounded corners=2pt,align=center,
              text width=0.265\linewidth,inner sep=6pt,font=\small},
  >={Latex[length=2mm]}, node distance=6mm]
\node[box] (limit) {\textbf{Finite limiting matrices}\\[2pt]
  Exact positivity on the\\four-particle space};
\node[box,right=of limit] (local) {\textbf{Fixed set of modes}\\[2pt]
  Exact change of variables;\\error bounded by pair energy};
\node[box,right=of local] (global) {\textbf{All particle sectors}\\[2pt]
  Add spectator modes and\\average over rotations};
\draw[->] (limit) -- (local);
\draw[->] (local) -- (global);
\end{tikzpicture}

\caption{The finite-sphere comparison preserves an error relative to the
pair energy at every step. The limiting positivity is imposed only after
wedging; no invertibility of its Gram matrices is required.}
\label{fig:transfer}
\end{figure}


Section~\ref{sec:algebra} derives the normal-ordered square, and
Section~\ref{sec:spin} establishes the spin formulas.
Section~\ref{sec:squares} gives the finite comparisons.
Section~\ref{sec:transfer} proves their uniform transfer, and
Section~\ref{sec:finish} controls the tail and proves the spherical results.
Section~\ref{sec:plane} derives the planar consequence.
The exact row data, rational formulas, and complete arithmetic procedure
for the positivity checks are in Appendix~\ref{app:certificate}.
